\chapter{Experimental Design}\label{ch:experimental-design}

This chapter fixes the configuration under which the method of
Chapter~\ref{ch:methodology} was instantiated: the corpus and its
split, the rosters, the conditions, the pre-registrations with what
each of them fixed in advance, and the statistical protocol.
Taxonomy-tuning mechanics and the construction ablation matrix are in
Appendix~\ref{app:tuning-protocol}; rubric-generation mechanics are in
Appendix~\ref{app:judge-generation}.

\section{Research Questions, Conditions, and Metrics}\label{sec:exp-rqmap}

Table~\ref{tab:rq-map} fixes which condition answers which research
question, with which metric, and which link of
Section~\ref{sec:research-questions} it lands on.

\begin{table}[htbp]
    \centering
    \caption{Research question to condition, metric, and link.}
    \label{tab:rq-map}
    \begin{tabular}{p{0.9cm}p{4.0cm}p{2.6cm}p{4.0cm}c}
        \toprule
        RQ & Question & Conditions & Primary metric & Link \\
        \midrule
        RQ1 & Do answer-key-blind, domain-conditioned pairwise
        judgments recover per-domain rankings, and what decides the
        verdicts?
        & C1 (main), with C5 (generic) as the rubric control
        & Per-verdict agreement with the MMLU-Pro key; Spearman $\rho$
        against ground-truth accuracy, reported under both tie
        conventions & 3 \\
        \addlinespace
        RQ2 & Does a cell-scoped rubric enable better judging
        than a generic rubric on the same questions?
        & C1 vs.\ C5, paired on an identical question set
        & Per-verdict key agreement; Spearman $\rho$ per domain
        reported alongside & 2 \\
        \bottomrule
    \end{tabular}
\end{table}

The scope is a \textbf{selected-domain study}: the arena evaluates
ex-ante-selected, high-support domains at high power rather than
all 14 at matched power. The selection rule is the four-criterion
substantive rule of Section~\ref{sec:exp-datasets}, applied inside a
fixed judge-call budget. Eight domains were run: mathematics, physics,
law, engineering, psychology, philosophy, history and computer science.
The set was fixed before the settled batch, under an offline
reordering screen computed at a since-superseded roster, and that
provenance is stated rather than laundered: refit at the settled
eight-model roster, the screen saturates --- the domain ordering equals
the global ordering ($\rho = 1.000$) in twelve of fourteen domains, and
only law ($\rho = 0.994$, $p = 0.003$) and chemistry ($\rho = 0.976$,
$p < 0.001$) show reordering at all, both marginal in magnitude. The
roster whose accuracy spread makes model pairs key-decidable also makes
domain membership nearly irrelevant to the ordering, so the domain
selection carries no endorsement from the refit screen and the eight
domains function as replications, not as screened contrasts
(Section~\ref{sec:res-law}). All eight ran against the same snapshot,
seed and roster under the paired protocol of
Section~\ref{sec:exp-arena-conditions}.

\paragraph{Run conditions, identical across all eight domains.}
Eight-model roster (Table~\ref{tab:roster-r8},
Appendix~\ref{app:models}), $\binom{8}{2} = 28$ model pairs;
\texttt{RESOLVED} option mode; the generic arm replays the cell-scoped
arm's exact $(\text{question}, \text{model pair})$ triples under the
same cells, so the arms share questions, comparisons and cells by
construction; the confidence gate disabled in both arms (no verdict was
removed by it in any domain); an isolated ranking database per run;
seed 42.

A decisive difference of $|\Delta\rho| \geq 0.007$ --- two steps of the
Spearman grid at the roster size it was written for --- was registered
in advance and is \textbf{withdrawn}: two runs of an identical
configuration, differing only in a scheduler seed, moved $\Delta\rho$
by six and zero grid steps, a spread three to four times the threshold
(Appendix~\ref{app:preregs}, addendum to P4). Throughout,
$\Delta\rho = \rho_{\textsc{main}} - \rho_{\textsc{c5}}$ is reported
for continuity and is not used as a decision rule; the reading
conventions that govern all of this are collected in one box at the
head of Chapter~\ref{ch:results}.

\section{Implementation Environment}\label{sec:exp-env}

The pipeline is implemented in Kotlin on Spring Boot, backed by SQLite
in WAL mode (Section~\ref{sec:arch-infra}). Embeddings are produced by
Qwen3-Embedding served locally; the judge is Mistral-Large-3, a single
large instruction-tuned model accessed through an Azure
OpenAI-compatible client, and the same judge model is used for rubric
induction and for every arena condition.\footnote{The judge identity
was verified against two independent sources after two project
documents were found to disagree: the Run~B run log records
\texttt{Mistral-Large-3} on every judge call (four occurrences of the
identifier, none of any smaller ``ministral'' variant), and the runtime
configuration pins \texttt{judge-model: \textquotedbl{}Mistral-Large-3\textquotedbl{}}. The roster
tables of Appendix~\ref{app:models} state it alongside the rosters.}

\textbf{Every construction and arena result in this thesis is at seed
42.} No result is reported as a mean over multiple seeds, and no
three-seed convention is in force. Seed dispersion has been measured
twice, and neither measurement is used as an error bar on a headline
number: a five-seed ablation inside the acceptance-gate analysis
(Section~\ref{sec:res-construction}), and a twenty-seed sweep of the
canonical configuration reported as the evidence for R3
(Section~\ref{sec:arch-requirements}).
The canonical configuration values are tabulated with their provenance
in Appendix~\ref{app:tuning-protocol}
(Table~\ref{tab:hyperparams}). Per-run identifiers, seeds, and realised
call counts are recorded in the manifest emitted with each run.

\section{Corpus, Split, and Information Separation}\label{sec:exp-datasets}

The corpus is MMLU-Pro~\parencite{wang2024mmlupro}: 12{,}032 questions
across 14 domains with human-verified reference answers, used in two
roles. It supplies the ground-truth domain labels that seed the 14
anchor nodes, and it supplies the per-domain accuracy scores used as
the post-hoc validation oracle (Section~\ref{sec:meth-validation}).
All 14 categories seed anchors, including the semantically
heterogeneous \textit{Other} category, which is retained as a geometric
anchor --- its queries remain routable mass and may migrate toward
coherent anchors --- but is not an arena selection candidate.
AG News~\parencite{zhang2015agnews} was used once, as a pipeline sanity
check on an independently labelled corpus; the check passed and the
corpus is not used further.

Taxonomy construction runs on all 14 domains simultaneously, which is
what allows queries to migrate between anchors and so makes the adapted
partition differ from the canonical one at all.
The arena evaluates a roster on ex-ante-selected domains,
chosen before observing any arena result under a fixed rule: high
held-out query support; high geometric coherence relative to siblings;
non-trivial adapted-versus-canonical migration volume; and
non-degenerate model performance variance under existing MMLU-Pro
accuracy. Math was the primary candidate under that rule. The offline
domain screen of Section~\ref{sec:res-law} was added later and is a
stronger instrument than criterion~(iv), because it asks directly
whether a domain's models reorder relative to their global ordering. It
flags four domains --- psychology, physics, history and law --- selects
law as the first to run, and demotes Math to a null arm. Its reach is
narrower than that use suggests, and the limit is registered rather
than discovered afterwards: the screen classifies a domain
\emph{aggregate} and says nothing about whether cells inside a domain
differ, the correlation between the two being $+0.072$. It is not
claimed as a clean predictor anywhere in this thesis
(Section~\ref{sec:res-leaf-null}).

\subsection{The Split}

The corpus is partitioned into two disjoint, stratified subsets:
\textbf{construction} (taxonomy induction and rubric generation) and
\textbf{arena-test} (held-out routing and judging only). The split is
$70/30$ --- $8{,}673$ construction against $3{,}445$ held out --- and is
stratified by MMLU-Pro top-level category: the corpus is shuffled
inside each of the 14 categories and cut there, so every category
contributes the same proportion to each side. It is not stratified any
finer than that, and in particular not by sub-topic tag.

\textbf{Tuning shares the construction side, and that is a
limitation.} A three-way split reserving a separate tuning-validation
subset would be cleaner, and it was rejected because at this corpus size
it starves the part that matters: the held-out pool already yields only
$131$ routable questions in the smallest domain, and taking a further
slice off it would leave several domains with too few questions to judge
at all. Instead, tuning gates score configurations on
construction-internal metrics, so no held-out query is seen by any gate.

The property the answer-key-blind claim rests on therefore holds, in
exactly the scope the code enforces. The reserved filter lives
in \emph{rubric induction}: no reserved question's text reaches the pool a
rubric is written from. It does \emph{not} exclude reserved material from the
partition, whose leaf regions do contain it. So the claim is that no
arena-test question influenced the rubric it is later judged under --- not
that the tree was drawn without reserved data.

That claim is verified by measurement rather than by the assertion alone: of
the $1{,}569$ distinct question texts the settled batch judged, \textbf{zero}
appear in any leaf's induction pool.
What is lost is the ability to distinguish a configuration that
generalises from one that suits the construction set: the tuning
decisions are validated on the same data that produced them, and a
reader should treat the chosen hyperparameters as fitted rather than as
independently confirmed.

One further detail affects what a ``construction-set'' claim means:
the rubric-induction pool is additionally thinned by a hash filter, so
a leaf's rubric samples its construction queries
(Section~\ref{sec:arch-judge}).\footnote{And about $4.2\%$ of MMLU-Pro
rows --- some 500 of 12{,}032 --- carry no link into the evaluation
tables; they are believed to be sentinel identifiers, but that has not
been confirmed, and they are recorded here as unattributed rather than
explained.}

The arena-test fraction targets approximately 10 held-out queries at
the smallest leaf under a selected domain. The binding constraint is
the number of model pairs per leaf, $\binom{8}{2} = 28$ at the settled
roster, and the bootstrap floor covers all 28 before adaptive
scheduling begins. The eleven-model figure of $\binom{11}{2} = 55$
applies only to the reliability-constant fit
(Section~\ref{sec:disc-debt}).

\subsection{Information Separation}\label{sec:exp-info-separation}

Table~\ref{tab:info-separation} makes the answer-key-blind design
concrete: at every pipeline phase, only the data marked \checkmark{} is
accessible to the corresponding component.

\begin{table}[htbp]
    \centering
    \caption{Data access by pipeline phase.}
    \label{tab:info-separation}
    \begin{tabular}{p{3.6cm}cccc}
        \toprule
        Phase &
        $\mathcal{Q}_{\text{con}}$ text &
        $\mathcal{Q}_{\text{con}}$ labels/answers &
        $\mathcal{Q}_{\text{eval}}$ text &
        $\mathcal{Q}_{\text{eval}}$ labels/answers \\
        \midrule
        Taxonomy construction & \checkmark & \checkmark & ---\textsuperscript{a} & --- \\
        Rubric generation & \checkmark & \checkmark & ---\textsuperscript{b} & --- \\
        Held-out routing & --- & --- & \checkmark & --- \\
        Arena evaluation (judging) & --- & --- & \checkmark & --- \\
        Post-hoc validation & --- & --- & \checkmark & \checkmark \\
        \bottomrule
    \end{tabular}
    \par\smallskip
    \footnotesize \textsuperscript{a}Construction is run on the construction
    split. The reserved \emph{filter}, however, is applied at rubric induction
    and not here, so this row holds by which pool construction was given rather
    than by an enforced exclusion. Because roughly $2{,}700$ question texts
    carry both a construction and a reserved identifier
    (Section~\ref{sec:exp-datasets}), a text that is also reserved under some other
    identifier can appear in a leaf region.
    \textsuperscript{b}Enforced, and verified: zero of the $1{,}569$ question
    texts the settled batch judged appear in any leaf's induction pool.
\end{table}

The evaluation set's labels and reference answers are accessed only in
the final, post-hoc validation phase, after every verdict has been
recorded. The rubric-induction filter that enforces the second row is
audited by a build-failing assertion that recomputes the membership
test independently and stops the build if the two disagree
(Section~\ref{sec:arch-judge}).

It guarantees exactly what it says and no more. The judge receives
neither the answer key nor the multiple-choice options --- its prompt
is the question stem, the two responses, and the output instructions.
What it does receive is responses that usually state their own answer,
$97.6\%$ of stored traces ending by declaring the answering model's
selected option (range $93.2$--$99.8\%$ across the eight domains). So ``answer-key-blind'' is not ``answer-blind'': the
key is withheld at the prompt and not at the mechanism, and
Chapter~\ref{ch:results} turns on the difference.

\section{Rosters and the Exclusion Policy}\label{sec:exp-rosters}

One roster carries every settled result: the eight models of
Table~\ref{tab:roster-r8} (Appendix~\ref{app:models}), selected to
span ground-truth accuracy $16.5\%$ to $90.5\%$ so that most model
pairs are key-decidable at the question counts a domain routes. What
that selection gives up is stated with the roster: it is not
length-matched (median stored response $220$ to $1{,}494$ characters,
$6.79\times$), so verbosity and capability are confounded in the
settled runs (Section~\ref{sec:disc-debt}).

Two earlier rosters existed and both belong to withdrawn runs. The
\textbf{pilot roster} (8 models) carried the capture-gap defect ---
four of its models had no stored response text, and the collinearity
that produced is tabulated with the historical rosters in
Appendix~\ref{app:models} (Table~\ref{tab:model-roster-pilot}). The
\textbf{length-matched band} (12 models) was assembled on the
repaired corpus to hold response length and accuracy close to
uncorrelated within the roster ($r = +0.286$, against $+0.597$ across
all 36 valid models); the campaign it ran is withdrawn, and the
length control it demonstrates is the one property of that roster the
settled selection sacrificed. All members of every roster produce
substantive responses; answer-only systems are out of scope for this
evaluation design, because a rubric has nothing to grade in a
response with no content.

\textbf{The eleven-model band} is a distinct, earlier roster still.
It was assembled at full law coverage, before an envelope unwrap
exposed one of its members as mixed-format, and it is the roster on
which the offline domain screen and the reliability constant were
first fitted (Appendix~\ref{app:models}). Every settled arena run
sits on the eight-model roster of Table~\ref{tab:roster-r8}.

That roster mismatch has since been resolved for the screen and not
for the constant. \textbf{The screen is refitted at the settled
eight-model roster} and saturates there
(Table~\ref{tab:domain-screen}); its earlier generations are
superseded and no screen $p$-value from them is carried forward.
\textbf{The reliability constant was not refitted.} $c = 2.52$
remains an eleven-model value, it is a property of its roster, and
the obligation to refit it stands (Section~\ref{sec:disc-debt}).

\paragraph{Ground-truth accuracy, and one accounting decision inside
it.} Every $\rho$ in this thesis is computed against per-model
per-domain accuracy exported from the upstream MMLU-Pro evaluation
outputs.\footnote{By \texttt{tools/analysis/export\_gt\_scores.py},
into \texttt{model\_domain\_scores.csv}.} A response whose predicted option
cannot be parsed is scored incorrect. That is not free: on the settled
roster, blank-prediction rates run $17.4\%$ to $19.7\%$ for its three
weakest models and below $1\%$ for the other five, so scoring blanks
as wrong deflates those three by up to $11$ accuracy points. Excluding
blanks from the denominator instead leaves the roster ordering
\emph{identical in all eight positions} on the reserved pool, which is
why every rank-based result here is unaffected by the choice. Both
columns are exported so the reader can check that claim.

\paragraph{Exclusion is enforced, and it is a measurement, not a
housekeeping decision.}
Ten models are excluded at roster load by a hard-coded list that
raises, each entry carrying its reason: one identifier-space mismatch,
six archives with no stored reasoning text, and three mixed-format
models. The mixed-format threshold is format \emph{consistency}, not
text presence: a model that emits a bare answer on a third of items
reintroduces the format preference inside its own comparisons. The
three excluded on that ground emit bare answers on $68.4\%$, $35.4\%$
and $27.6\%$ of items respectively.

The honest statement of the policy is therefore:
\emph{response text was not persisted by the generation pipeline for
several models, and three further models emit a bare answer on more
than a quarter of items; all are excluded, and the format artifact is
reported as the measurement that justifies the exclusion.}

One note belongs with it: six of the ten exclusions are for exactly
the condition the corpus validator's trace-presence check measures and
declines to act on, because that check has no failure branch. The
policy is a ratchet on known failures, not a threshold, and it does
not generalise to any new model.\footnote{And one disambiguation: the
roster member \texttt{Meta-Llama-3\_1-70B-Instruct} is not the
excluded \texttt{Meta-Llama-3-70B-Instruct}. They are different
models, the roster's is clean, and the excluded one sits in a
different question-identifier space and is what collapsed law to 20
questions in Section~\ref{sec:res-law}.}
The full roster tables are in Appendix~\ref{app:models}.

\section{Conditions}\label{sec:exp-arena-conditions}

Three conditions appear in this thesis: C1 (\textsc{main}), C3
(\textsc{generic}), and C5, registered as \textsc{no-partition} and
more accurately described as a generic-rubric arm on the same partition
(see below). Every paired run uses two of them, C1 against C5. C3 ran once, on the pilot. Each chapter names them in full at first use
and by the small-caps short form thereafter; run exports spell the
generic condition \texttt{GENERIC\_JUDGE}, and the two names denote
the same condition.\footnote{The identifiers are inherited from an
earlier design document and are deliberately not renumbered: C2 ---
re-aggregating C1's verdict set by editorial label --- and C4, an
optional diagnostic, were cut before any run. C2 could not have
isolated anything: it shares its entire verdict set with C1, and the
question it was built to answer --- the per-cell
adapted-versus-canonical rank difference --- was refuted at its
premise under pre-registration before the arena ran
(Section~\ref{sec:res-granularity}).}
Table~\ref{tab:arena-conditions} states what each condition varies.

\begin{table}[htbp]
    \centering
    \caption{Arena conditions. All conditions share the frozen
    snapshot and the model roster of their run; only the listed factor
    varies.}
    \label{tab:arena-conditions}
    \begin{tabular}{p{0.9cm}p{2.7cm}p{4.4cm}p{4.0cm}}
        \toprule
        ID & Condition & Delta from C1 & What it controls \\
        \midrule
        C1 & Main
        & --- (adapted partition, cell-scoped reference-informed rubric)
        & reference \\
        \addlinespace
        C5 & Generic rubric
        & Replays \textsc{main}'s exact triples under the same cells,
        judged under a generic MT-Bench-style pairwise prompt
        & rubric bundle at matched budget; \textbf{not} the partition \\
        \bottomrule
    \end{tabular}
\end{table}

A third condition, C3 (a rubric-free judge on the same cells), ran
only inside the withdrawn pilot and carries no result; the isolation
run that would revive its idea on working code is named as future
work (Section~\ref{sec:disc-future}). What C5 controls is stated in
the table in its corrected form: under the replay design of
Section~\ref{sec:exp-rqmap}, pooling per-cell sufficient statistics is
arithmetically identical to a direct fit at the domain, so C5 varies
the judge-prompt bundle on a fixed partition and nothing else. No
condition in this thesis removes the partition: a real no-partition
arm would schedule its own comparisons over the whole domain with no
cells, holding the rubric fixed --- a new scheduler condition, not a
configuration change, and it was not run.
Figure~\ref{fig:design-schematic} draws the design end to end; the
one fact it exists to make unmissable is that the two arms read a
single comparison list.

\begin{figure}[htbp]
\centering
\footnotesize
\begin{tikzpicture}[
  box/.style={draw, rounded corners=1pt, align=center,
              inner sep=4pt, font=\footnotesize},
  lbl/.style={font=\scriptsize, text=black!60},
  arr/.style={->, >=stealth, thick},
  node distance=4mm and 9mm]
\node[box] (corpus) {MMLU-Pro with eval coverage\\$12{,}118$ questions};
\node[box, below left=8mm and -14mm of corpus] (con)
  {construction split\\$8{,}673$ ($70\%$)};
\node[box, below right=8mm and -14mm of corpus] (res)
  {reserved split\\$3{,}445$ ($30\%$)};
\node[box, below=of con] (tree)
  {frozen taxonomy + rubrics\\$87$ cells, snapshot \texttt{20260727}};
\node[box, below=of res] (routed)
  {routed held-out questions\\$1{,}637$ in $52$ cells};
\node[box, below=14mm of $(tree.south)!0.5!(routed.south)$] (triples)
  {\textbf{one comparison list}\\$5{,}161$ (cell, pair, question) triples};
\node[box, below left=8mm and -16mm of triples] (main)
  {\textsc{main}\\cell rubric, JSON verdict};
\node[box, below right=8mm and -16mm of triples] (gen)
  {\textsc{generic} (C5)\\MT-Bench prompt, \texttt{[[A]]/[[B]]/[[C]]}};
\draw[arr] (corpus) -- (con);
\draw[arr] (corpus) -- (res);
\draw[arr] (con) -- node[lbl, left] {induction (+ tuning)} (tree);
\draw[arr] (res) -- node[lbl, right] {routing, by the frozen geometry} (routed);
\draw[arr] (tree) -- (triples);
\draw[arr] (routed) -- (triples);
\draw[arr] (triples) -- node[lbl, left, pos=0.6] {judged} (main);
\draw[arr] (triples) -- node[lbl, right, pos=0.6] {replayed} (gen);
\node[lbl, below=2mm of $(main.south)!0.5!(gen.south)$, align=center]
  {the arms share cells, comparisons and aggregation;\\
   they differ in the judge-prompt bundle and in nothing else};
\end{tikzpicture}
\caption{The experimental design. Rubric induction reads only the
construction split (verified: zero of the $1{,}569$ judged texts
appear in any induction pool); held-out questions are routed by the
frozen geometry; and both judging arms read \emph{one} comparison
list, so the contrast isolates the judge-prompt bundle at matched
budget. No condition removes the partition.}
\label{fig:design-schematic}
\end{figure}

\paragraph{Pairing, and what may be compared across formats.}
Both arms of the C1-versus-C5 comparison judge an identical question
set --- the replay design of Section~\ref{sec:exp-rqmap}. Unlike the
original MT-Bench protocol, which samples independently, here the
rubric bundle is the treatment and the partition is shared
infrastructure, so replay is what isolates the treatment; independent
sampling would add variance orthogonal to the thing under test.
The two arms parse different verdict formats --- a structured schema
against \texttt{[[A]]}/\texttt{[[B]]}/\texttt{[[C]]} markers --- with
three consequences. Winners, the Bradley--Terry ordering derived from
them, and tie counts \emph{are} comparable, provided ties are defined
identically on both sides. Judge confidence is \emph{not}: the generic
arm emits a flat constant. Parse-failure rates are not comparable
either, and the failure \emph{modes} differ rather than merely the
rates, because a malformed structured response fails loudly while a
generic response with no marker returns an invalid verdict silently.
Operationally this forces one thing that is easy to miss: the
confidence gate must be disabled in both arms, since at a flat
confidence it would admit every generic verdict while filtering the
main arm's, and the arms would no longer be paired where it counts.

\section{Pre-Registration}\label{sec:exp-prereg}

Five decision rules were fixed in dated documents before the data they
govern existed, and four of them govern runs that happened. The dates
are part of the claim and are stated from the version-control history:
the granularity analysis, the rubric-specificity test, and the arena
launch-cost bands were registered on 2026-07-27; the no-partition
baseline was registered on 2026-07-27, the same day the domain screen
settled on law (Section~\ref{sec:res-law}), and amended through
2026-07-30. The documents are reproduced in
Appendix~\ref{app:tuning-protocol}; what each fixed is summarised here,
because Chapter~\ref{ch:results} calls results ``pre-registered'' and a
reader is entitled to check the claim.

Five dated addenda extend the rubric-specificity baseline. The first
is the results-written-back addendum to P2 (2026-07-28), described
with P2 in Appendix~\ref{app:preregs}. Of the four that extend P4,
each but the last was written before the runs it governs: the second
(2026-07-29) registers
psychology as a reordering domain and engineering as a null. The third
(2026-07-30) registers physics as the last untested flagged domain and
computer science as a third null, and refits the screen at the
then-current twelve-model roster --- a refit itself since superseded
by the settled eight-model one of Section~\ref{sec:res-law}. The fourth (2026-07-30) registers the mathematics re-run
against the coverage confound described below, and fixes in advance how
each of its three possible outcomes is to be read. The fifth
(2026-07-30) is written after seven runs have results; it records the
outcomes against the predictions, retires chemistry as registered but
never run, and states explicitly that philosophy and history were never
registered prospectively and are reported as replications.

\textbf{Granularity and discriminative power.} Three outcomes were
named in advance, and the primary statistic was fixed as the
\emph{observed} correlation rather than the disattenuated one. That
rule is the only reason the conclusion is readable: the disattenuated
column moves in the opposite direction and clips past $1.0$, so leading
with it would have produced a conclusion from the arithmetic of the
correction rather than from the data.

\textbf{Rubric specificity.} Four directional predictions, with the
confound direction named \emph{per measure} --- including one measure
whose vocabulary-size confound could have driven the null arm
\emph{higher} --- plus within-arm decision thresholds.

\textbf{Arena launch cost.} Cost bands with the instruction that no
band is to be renegotiated after the number is seen, and a void
condition: if the measured per-cell identification cost came back
constant across cells, the scheduler would still be capping and the
measurement would report nothing.

\textbf{The rubric-specificity baseline.} The paired design; the primary
and secondary outcomes; the registered decisive difference of
$|\Delta\rho| \geq 0.007$, since withdrawn (Appendix~\ref{app:preregs},
addendum to P4); four void conditions (identical question sets,
confidence gate off in both arms, identical samplers, identically defined
ties); and, recorded before the run, the honest prior that with options
present the comparison would likely come back null, together with the
reason.
Two of this registration's provisions are withdrawn in the dated
addendum (Section~\ref{sec:prereg-addendum-p4}); the registered text
above is preserved unedited, as a registration must be.

\textbf{The mathematics coverage confound, and the re-run registered
against it.} In the withdrawn twelve-model campaign, mathematics was
the only run predating the two scheduler repairs --- the mandatory
bootstrap coverage floor with its in-run assertion, and per-(leaf,
pair) query shuffling;\footnote{Commits \texttt{5e1be09} and
\texttt{a5c36fd} respectively.} its comparison graph was a leader-centric star covering
under a third of the model pairs, and no property of mathematics
could be inferred from it. That campaign is withdrawn in its
entirety, so the confound is recorded here as the reason the re-run
registration exists rather than as a qualification on any reported
number. The fourth addendum
registers a re-run on the repaired scheduler under the original null
prediction and fixes three readings in advance: a reproduced negative
means the failed null is real and now rests on a complete graph; a null
means the original was a scheduler artifact and is superseded on stated
grounds; a positive means the screen failed in the opposite direction
from philosophy. Reporting only the re-run is excluded under every
outcome. That run was outstanding when this chapter was fixed.

The baseline conditions for taxonomy quality specified in an earlier
design --- flat $k$-means, agglomerative clustering, and a random
taxonomy null, scored by a family of reference-free hierarchy metrics
--- are not reported. The supervised hierarchical-classification
baselines were removed as unimplemented; no gold taxonomy exists for
MMLU-Pro, so reference-dependent metrics were dropped rather than
computed against an ad-hoc structure; and the remaining runs were never
executed. None of those metrics bears on any of the three links of
Section~\ref{sec:research-questions}, and the construction evidence
reported instead is the fixed-point certificate, the leaf-size
distribution, and the separation nulls of
Section~\ref{sec:arch-requirements}.

One housekeeping note on the rubric-specificity document. As
originally registered it recorded ``treatment arm only, no null, this
is not yet evidence''; the null arm was subsequently run, and a dated
results-written-back addendum (2026-07-28) now closes the document:
the headline separation ($p = 1.21 \times 10^{-6}$) was reproduced
from the stored artifacts on 2026-07-27, and the addendum records that
those null-arm artifacts sit in an untracked build directory --- a
reproducibility caveat the document carries explicitly
(Section~\ref{sec:res-rubric-null}).

\section{Metrics and Statistical Protocol}\label{sec:exp-metrics}

\textbf{The primary judge-side metric is per-verdict agreement with the
MMLU-Pro answer key}, on the subset where exactly one response is
correct and the judge did not tie. It is a proportion over roughly a
thousand verdicts per condition. At $p \approx 0.85$ and $n = 1000$ the
per-arm standard error is about $1.1\%$, so the standard error of the
paired difference between two arms is about $1.6\%$. Observed gaps run
from $0.9$ to $4.0$ points, which is why this metric is described as
leaning rather than as deciding wherever it is quoted.

\textbf{Spearman $\rho$ against ground-truth accuracy is the partition
metric, and it is heavily quantised.} At $M$ models the grid spacing is
$6/(M(M^2-1))$: $0.0119$ at $M = 8$, $0.0046$ at $M = 11$, $0.0035$ at
$M = 12$. A single Spearman value therefore consumes a full arena run to
produce one number resolvable only to an adjacent transposition, which
is why it is the secondary statistic everywhere except the partition
claim, where it is the thing being claimed about.

% RESOLVED 2026-07-31, and not in the way this note anticipated. The seam
% between the literal 0.007 and two grid steps (0.0069930) no longer needs a
% ruling, because the threshold itself is withdrawn: two runs of an identical
% configuration differing only in a scheduler seed gave Delta rho of six and
% zero grid steps. A rule at two steps cannot adjudicate a quantity whose
% run-to-run spread is three to four times larger. The paragraph below now
% states that rather than the ambiguity.

\textbf{The threshold is withdrawn, and the reason supersedes the
ambiguity it had.} The pre-registration expressed the decisive difference
two ways at once --- as the number $0.007$ and as two steps of the
Spearman grid --- which differ by seven millionths and which several
contrasts landed exactly between. That seam no longer needs resolving,
because the threshold did not survive the seed replicate of
Section~\ref{sec:exp-rqmap}: a rule set at two grid steps cannot
separate signal from run-to-run noise three to four times its size.
$\rho$ is therefore reported for continuity with prior work and is not
used as a decision rule.
The replicate behind that statement: two runs of the history domain at
an identical configuration in \texttt{RESOLVED} mode, seeds 42 and 137,
moved $\Delta\rho$ by six and zero grid steps while per-verdict
agreement reproduced to within $0.1$ percentage points.

\textbf{Bradley--Terry $\theta$, never win-rate}, for the
scheduler-bias reason stated once in
Section~\ref{sec:meth-scheduling}; the measured cost of the difference
is reported beside the result it affects
(Section~\ref{sec:res-c5}).

\textbf{Every $\rho$ is reported under both tie conventions}, per
rule~D5 of Section~\ref{sec:meth-discipline}; the measured instability
that forces the rule is Section~\ref{sec:res-tie-policy}.

\textbf{No confidence interval is reported anywhere.} The
Bradley--Terry variance computation subtracted $1/K$ where the rank
completion requires $1/K^{2}$, so no standard error computed before the
fix is a measurement, and the intervals built on those standard errors
were never recomputed after it (Section~\ref{sec:meth-bt}). Standard
errors on proportions, which are binomial and not Bradley--Terry, are
reported and are labelled as such. Every bracketed range in this thesis
is an interquartile range and says so.

\textbf{Disattenuation is two-sided.} Disattenuation corrects a
correlation for the measurement error in the two rankings it compares.
Both sides of
$\rho(\text{arena ranking}, \text{ground-truth ranking})$ are
finite-sample per-cell rankings, so the correction is $\rho_{\text{obs}}
/ r$ with $r = n/(n+c)$, not $\rho_{\text{obs}} / \sqrt{r}$; the
one-sided form applies only when one measure is perfectly reliable. The
corrected constant is $c = 2.52$ at the 11-model band, and three
caveats travel with it wherever it is used: it is a property of the
roster and must be re-fitted if the roster changes; the interval
originally quoted for it resamples three fitted points and is not a
confidence interval; and it was fitted on ground-truth rankings, which
are noise-free per question, so it is a \emph{ceiling} on arena
reliability rather than a prediction of it
(Section~\ref{sec:disc-debt}).

\textbf{Bootstrap resampling counts differ by path and are stated where
used.} The validation path uses 2{,}000 resamples; the ranking path
uses 50; the in-run trajectory uses 1. A blanket claim of 2{,}000
resamples everywhere would be wrong, and figures derived from the
ranking path are quoted with their own count.

\textbf{No pass/fail threshold on $\rho$.} An earlier design
pre-specified $\rho > 0.7$ as a criterion. It is retained only as an
interpretive marker of strong agreement; results below it are reported
in full and scoped to the domains they were measured on.
